\chapter{Conclusion}

% ============================================================
% CONTENT MAP. NOTE: an earlier version of this map pointed to
% THESIS_MASTER_RECORD.md "Sec.12" for limitations -- CORRECTED:
% Limitations is Sec.11. Sec.12 is a SEPARATE "Open Items / Remaining
% Work" list (genuinely unfinished engineering, not limitations of
% completed work) -- use both, don't conflate them.
%
%   \section{Summary of Contributions}
%     - the three-attack, four-model, backend-controlled benchmark
%       itself (THESIS_MASTER_RECORD.md Sec.1-2)
%     - the backend-confound discovery as a citable methodological
%       warning (Sec.8, PHASE4_EU_AI_ACT_MAPPING.md Sec.3.4's "the
%       inference stack itself must be held constant, or a compliance
%       test could certify a defense that does nothing")
%     - the EU AI Act Article 15 technical mapping (PLAN.md rule 5 --
%       present as a technical contribution, not a compliance claim).
%       Reproduce PHASE4_EU_AI_ACT_MAPPING.md Sec.4's FOUR concrete
%       contributions precisely, don't paraphrase into something vaguer:
%         1. a demonstrated three-family RAG threat taxonomy that maps
%            onto Article 15(5)'s illustrative list while also
%            surfacing where that list's classical-ML framing doesn't
%            cleanly cover a RAG-relevant threat (Crescendo)
%         2. a demonstrated statistical testing discipline (paired exact
%            tests, Holm-Bonferroni correction) that could inform what
%            "adequate testing" means for Article 15(2)'s benchmark
%            mandate, in a literature that inconsistently reports
%            significance testing at all (per this project's own
%            CITATIONS.md survey)
%         3. a citable warning about the inference-backend confound as
%            a specific class of testing error a compliance regime
%            could walk into
%         4. an honest demonstration of what current, real,
%            off-the-shelf inference-time mitigations actually achieve
%            -- one structurally wrong for its threat class
%            (output_filter/PoisonedRAG), one costing severe utility
%            for a real effect (spotlighting), one mechanism-confirmed
%            but statistically underpowered (instruction_detection)
%       AND the explicit "what this project does NOT claim" paragraph
%       (Sec.4's own wording): not a legal compliance assessment, not a
%       conformity assessment under Article 43, does not establish any
%       deployed system's Annex III classification, does not test every
%       15(5) vulnerability class (model poisoning, confidentiality
%       attacks are named gaps) -- keep this paragraph, don't cut it for
%       length; it is what keeps the mapping honest (PLAN.md rule 5)
%
%   \section{Practical Implications}
%     -- reproduce THESIS_MASTER_RECORD.md Sec.10's own scope-note
%     framing: every claim here is "this project's results show X for
%     the models/attacks/defenses tested," never a general
%     recommendation. Source: Sec.10.1-10.4:
%     - 10.1: no single best model -- robustness ranking INVERTS
%       depending on which attack: llama-3.1-8b is the strongest
%       performer against injection and the WEAKEST against Crescendo
%       (70.2\%, highest of the four)
%     - 10.2: defense scorecard (reuse the precise per-defense numbers
%       already written into the Evaluation chapter map -- don't
%       re-derive a softer version here)
%     - 10.3: efficiency was never measured -- named explicitly as a
%       scope this project did not evaluate (latency, throughput,
%       memory), not folded into the robustness conclusions
%     - 10.4: don't assume one defense generalizes across attack types;
%       match the defense to the actual threat model; test at realistic
%       scale before trusting small-sample results; the EU AI Act
%       compliance-testing gap is current and unresolved
%
%   \section{Limitations}
%     -- THESIS_MASTER_RECORD.md Sec.11 (NOT Sec.12), full honest list,
%     not softened. Reproduce these precisely rather than summarizing
%     into vague hedges:
%     - nq_open exclusion still pending explicit supervisor
%       acknowledgment as of the source document -- confirm current
%       status with your supervisor before final submission and update
%       this framing if it's since been formally signed off
%     - the n asymmetry within injection (40 vs. 1000) -- a real
%       compute constraint, not a design choice
%     - PoisonedRAG's matched-subset reduction (n=90-96 baseline vs.
%       n=14-18 matched defended subset per cell)
%     - Crescendo's Phase 3 behavior-pool mismatch -- only 9 of 20
%       Phase 3 behaviors per model appear in the Phase 2 baseline,
%       usable paired n is 9 not 20
%     - ASR scoring is rule-based substring matching for all three
%       attacks -- confirmed to conflate genuinely different phenomena
%       under one binary flag
%     - is\_refusal()'s false-positive rate confirmed for
%       ministral-3-8b only, individually audited; plausibly affects an
%       unquantified fraction of the other 3 models' flagged rows
%     - one attacker/judge model (DeepSeek V4 Pro) for Crescendo, one
%       poison-generator model for PoisonedRAG -- cannot distinguish
%       "these models resist similarly" from "this specific
%       attacker/generator happens to land similarly"
%
%   \section{Future Work}
%     -- THESIS_MASTER_RECORD.md Sec.12 (Open Items), separate from
%     Limitations above:
%     - output_filter-for-Crescendo: built, guard verdict logged, but
%       deliberately never wired to enforce live in the conversation --
%       the clearest, most concrete next step (this is unfinished
%       engineering, not a limitation of a completed test)
%     - a robustness-efficiency trade-off study (Sec.10.3) -- named as
%       a well-motivated direction, not attempted here
%     - model poisoning and confidentiality attacks (Article 15(5)'s
%       other named vulnerability classes) as a natural extension of
%       this project's threat taxonomy
%
% Suggested word target: ~1,500 words
% ============================================================
